\section{Background}

\subsection{Lag-1 Online Recalibration and Its Hazard}

the prior online-calibration evaluation~\cite{paper7online} evaluated lag-1 rolling-history grid-search
recalibration for the HygienePrio scorer. At moderate capacity
($K \in \{50, 100\}$), the procedure recovered essentially all of the
offline-peek calibration ceiling (per-cell recovery ratio $\rho \in
\{1.04, 0.99\}$). At high capacity ($K = 200$), it produced a net
\textit{harm}: $-5.9$~pp at W2, $-7.8$~pp at W3, and a cell-mean
recovery ratio of $-0.66$. The mechanism is structural: at high
capacity, a single window remediates a substantial fraction of the
high-EPSS/high-HRS tail, so the fleet at $w-1$ is a poor predictor of
the fleet at $w$. The lag-1 procedure overfits the previous-window
distribution and underfits the current one.

\subsection{Smoothing as a Variance-Bias Trade-Off}

A standard fix for ``too-noisy single-window estimates'' is to average
across a longer window of history. Two operationally simple variants
appear in time-series forecasting~\cite{huang2024ewma}:

\textit{Trailing arithmetic mean.} Equal weight on the past $H$
windows; reduces variance by $1/H$ but introduces lag of order $H/2$.

\textit{Exponentially-weighted moving average (EWMA).} Geometric
decay $\alpha^i$ over past windows; trades smoothness for a tunable
effective horizon. At $\alpha = 0.6$ the effective lookback is
$\approx 2.5$ windows.

The question this paper asks is whether either smoothing scheme
reverses the prior online-calibration evaluation's high-capacity hazard \textit{without} breaking the
moderate-capacity recovery. A naive prediction is that smoothing
helps at high $K$ (less overfitting to noisy single-window targets)
and hurts at low $K$ (introduces unnecessary lag when single-window
calibration is already accurate); whether this trade-off resolves
favourably is an empirical question.

\subsection{Why This Matters Operationally}

An operations team that adopts the prior online-calibration evaluation's online recalibration as
written and runs it on a high-capacity fleet would suffer a measurable
P@50 penalty in the first weeks of operation. Any deployable
calibration recipe must therefore include a smoothing knob; this
paper evaluates the simplest two such knobs in a pre-registered
setting.
